\section{Completing tuples and rounding their phases}
\label{sec:rounding}

Return to the sequence of germs, with the tier decomposition of
Proposition~\ref{prop:product}. Degrees within a tier have uniformly
bounded ratios; between consecutive tiers their ratios tend to
infinity. Let $m$ be the end of a tier and $m_-$ the end of the
preceding tier, with $m_-=0$ for the first one. Write $d=d_m$.
Consider a tuple and an independent positive weight assignment as
in Proposition~\ref{val:optimization}, satisfying
\begin{equation}\label{round:input-bounds}
 \max\{1,m_-\}\le r\le m,\qquad
 f_1=f_1^0,\qquad \frac12\le\frac{b_1}{d_1}\le2,
 \qquad b_j\le B_*d_j.
\end{equation}
Assume it has a quasi-monomial Gauss extension of ratio at most
$M_*$. The bounds $B_*,M_*$ are prescribed before this application.
For its selected optimal vertex, \eqref{val:comparison} gives
\begin{equation}\label{round:s-bounds}
 c\le s(w)\le2,\qquad A_\eta(w)\le2M_*,\qquad
 c=\frac1{4nM_*}.
\end{equation}
Indeed the upper bound for $s$ follows by testing $f_1$, while
$b_1\le2nA_\eta(w)d_1\le2nM_*s(w)d_1$ gives the lower bound.

For reference, the notation and the stages of its use are as follows.
\begin{center}
\small
\begin{tabular}{@{}p{0.18\textwidth}p{0.78\textwidth}@{}}
\toprule
Notation & Role and scope \\
\midrule
$v,f^0,d_i$ & Reference valuation and degree scales from
\eqref{val:degree-data}; the stated lower, first-degree and product
bounds do not bound every $d_i$ above. \\
$f,b,w,b^*$ & Independent input weights and an optimal Gauss extension
at selection; the exact target $b^*$ may be dependent and is reached
by continuation, as in \eqref{round:targets}. \\
$s,\Phi$ & The comparison factor \eqref{val:comparison} and the
ratio \eqref{val:ratio}, both depending on the valuation. \\
$m_h,e,g_s$ & Rational coefficient vectors at a selected branch,
defined by \eqref{val:coefficient-vectors}; their denominator bound
there belongs to that fixed vertex. \\
$B_*,M_*$; $S,D_{\rm fin},\ell$ & Prescribed numerical budgets,
then the common rounding and clearing integers chosen below;
see \eqref{round:final-count-constant} and
\eqref{round:phase-valuation}. \\
\bottomrule
\end{tabular}
\end{center}
The denominator of a fixed selected vertex and the required reduction
characteristic may depend on the fixed germ, finite tests and cutoffs;
no uniform bound on their complexity or number of Newton breakpoints
is asserted. The budgets and the later $S,D_{\rm fin},\ell$ have
exactly the uniformity along the retained subsequence proved below,
whereas the final cutoff $T$ may depend on the function and germ.

In the following estimates, constants depend only on the prescribed
bounds, dimension, tier comparability, and the constants in
Proposition~\ref{prop:cone-estimates} and Lemma~\ref{val:window}.
For every fixed cutoff multiplier $K\ge1$ occurring below, omit
a finite initial part of the sequence so that $Kd$ is below the
next tier. With
\begin{equation}\label{round:Dj}
 \mathcal D_j=\prod_{i\le j}\frac d{d_i},
\end{equation}
Lemma~\ref{val:window} and comparability within the current tier give
\begin{equation}\label{round:window-tier}
 \mathcal H(Kd)\le C_WK^m\mathcal D_m,\qquad
 \mathcal D_m\le C_2\mathcal D_r.
\end{equation}
When $m=n$, no restriction from a next tier is needed.

\subsection{Three uniform counting estimates}

Suppose the tests in Proposition~\ref{car:output} have powers with
\begin{equation}\label{round:test-cost}
 \sum_\nu\frac{j_\nu}{p}w(t_\nu)\le d.
\end{equation}
For sufficiently large reductions its output satisfies
\begin{equation}\label{round:R-cost}
 w(R)\le(2+rB_*)d.
\end{equation}
In fact \eqref{car:C1} and \eqref{car:C2-value} have leading
contribution $(\phi g_s-e)\cdot b=0$; the fixed-data error divided
by $p$ tends to zero, and $0\le k_i/p<1$. The box of functions
\begin{equation}\label{round:basic-box}
 Rf^\alpha,\qquad
 0\le\alpha_i\le d/d_i,\quad\alpha_i\in\Z,
\end{equation}
has weight at most $(2+2rB_*)d$. Fix a uniform $K$ large enough
that this is at most $s(w)Kd$, using \eqref{round:s-bounds}.
There are at least $\mathcal D_r$ distinct exponent-character
pairs in the box, since $b$ is independent. Each product is a
regular eigenfunction. The cutoff inclusion
\eqref{car:cutoff-inclusion} shows that products with distinct
(value, character) pairs have linearly independent classes modulo
$I_{Kd,p}$: split a putative dependence into characters; within one
character its unique lowest value is at most $sKd$, contrary to
the inclusion. Thus their number
is at most $\mathcal H(Kd)$. By \eqref{round:window-tier}, at most
a uniform number $C_3$ of such boxes can be pairwise disjoint in
their exponent-character pairs. All cutoffs invoked here are fixed
before taking the reduction characteristic large.

\begin{lemma}\label{round:algebraic-counts}
For regular eigenfunctions $t$ whose $w$-initial is algebraic over
the graded tuple field of $f$, put
\[
 \Lambda_f=\left\langle(\mathbf e_i,\chi_{f_i}):1\le i\le r
                     \right\rangle_{\Z}
 \subset\Q^r\times\operatorname{Hom}(G,k_0^*),
\]
where $\mathbf e_i=m_{f_i}$ is the $i$th standard vector. Then
\begin{align}
 \#\{(m_t,\chi_t)\bmod\Lambda_f\}&\le C_3,
 \label{round:U1}\\
 |m_{t,i}|&\le C_4\frac{w(t)}{d_i}\quad(1\le i\le r).
 \label{round:U2}
\end{align}
\end{lemma}

\begin{proof}
For \eqref{round:U1}, take any finite list of functions representing
distinct input classes. Use them as the tests in
Proposition~\ref{car:output}. For each test take its power one and
all other powers zero, using the same $H$ for these outputs.
For arbitrarily large $p$, \eqref{round:test-cost} holds. If two
output classes were equal, subtracting
\eqref{car:C2-value} and \eqref{car:C2-character} would show that
the corresponding input classes are equal: the $k$ terms belong
to $\Lambda_f$, and multiplying an element of $\Lambda_f$ by
$p$ remains in $\Lambda_f$. No division by $p$ in a character
group is used. Thus these outputs give disjoint boxes
\eqref{round:basic-box}, and the size of every finite list is
at most $C_3$.

If $w(t)=0$, independence of $b$ gives $m_t=0$. Otherwise apply
the same proposition to the powers of this one test with
$0\le j/p\le d/w(t)$. Its output's $i$th exponent coordinate
is a common translate plus $(j/p)m_{t,i}$ and an error in
$[0,1)$. The common $o(1)$ error from \eqref{car:C1} is independent
of $j$. If $|m_{t,i}|d/w(t)$ exceeded a sufficiently large fixed
multiple of $d/d_i$, choose more than $C_3$ evenly spaced
values of $j/p$ in this interval. For large $p$ they can be
rounded to the $1/p$ grid with arbitrarily small additional error.
Since $d/d_i\ge1$, their output coordinates can be made more
than $d/d_i$ apart. The corresponding boxes are disjoint,
contradicting their uniform bound. This proves \eqref{round:U2}.
\end{proof}

\begin{lemma}\label{round:enlarge}
If $r<m$, there is a centred regular eigenfunction $t$ with
$w(t)\le C_5d$ whose initial is transcendental over the graded
tuple field. Adjoining $t$ preserves the Gauss property at $w$.
The enlarged tuple can then be continued to arbitrarily close
independent assignments, with ratio increased by at most one.
\end{lemma}

\begin{proof}
Suppose all regular eigenfunctions of $w$-weight at most $L_*d$
have algebraic initial. At any fixed pair $(m_t,\chi_t)$ the
initials span at most one dimension: after adjoining $x$ and
normalizing, \eqref{car:constant-tests} puts them in $k_0$.
By \eqref{round:U1} there are at most $C_3$ classes, and within
one class the shift in $\Lambda_f$ is determined by its integral
exponent vector. By \eqref{round:U2}, the number of possible
pairs is at most
\begin{equation}\label{round:upper-count}
 C(1+L_*)^r\mathcal D_r.
\end{equation}

Now use the original tuple, not the changing optimized one.
For a fixed large $D\ge1$, consider the span of
\[
 (f_1^0)^{\alpha_1}\cdots(f_m^0)^{\alpha_m},
 \qquad 0\le\alpha_j\le Dd/d_j.
\]
Its dimension is at least $D^m\mathcal D_m$. Every nonzero
combination has $v$-value at most $mDd$, because the original
initials are independent. By \eqref{val:comparison} and
\eqref{round:s-bounds}, its $w$-value is at most $C'Dd$.
Decompose this finite-dimensional space into character subspaces
and filter each separately by $w$. This does not require the
$w$-filtration itself to be group-invariant. A valuation
filtration on a finite-dimensional vector space has a basis
adapted to its finitely many occurring values, so the dimensions
of its initial spaces sum to its dimension. If all these
eigenfunction initials were algebraic, \eqref{round:upper-count}
with $L_*=C'D$ would bound that dimension by
$C(1+C'D)^r\mathcal D_r$. Since $r<m$ and
$\mathcal D_m\ge\mathcal D_r$, a sufficiently large fixed $D$
makes this impossible. This supplies $t$ with the asserted
uniform bound.

The function is a nonunit, because the initial of a unit centred
at the closed point is a nonzero scalar. Its transcendental
initial enlarges the Gauss tuple. The degree $d_{r+1}$ belongs
to the current tier, so $d/d_{r+1}$ is uniformly bounded.
Finally apply Lemma~\ref{val:gauss-tools}, retaining the $h_i$,
to move to independent weights as close as desired. Discrepancy
and $s>0$ are continuous in that continuation, so one can
reserve at most one unit of increase in $\Phi$.
\end{proof}

\begin{lemma}\label{round:derivative}
When $r=m$, the selected vertex satisfies
\begin{equation}\label{round:U3}
 \left|\frac{(e-\phi g_s)_i}{s}\right|
       \le C_6\frac d{d_i}\quad(1\le i\le m).
\end{equation}
Consequently, on a selected ratio branch along a segment of
weights, at its independent parameters,
\begin{equation}\label{round:derivative-bound}
 \left|\frac{d\Phi}{dt}\right|
 \le C_6\sum_{i\le m}\frac d{d_i}|\dot b_i|.
\end{equation}
\end{lemma}

\begin{proof}
Apply Proposition~\ref{car:output} with no $t_\nu$. Its output
has
\[
 m_R=\phi g_s-e+k/p+o(1),\qquad
 w(R)\le(2+rB_*)d.
\]
Fix a sufficiently large integer $D$ and consider, in the
reduction, all boxes
\[
 R^jf^\alpha,\qquad 0\le j\le D,\qquad
 0\le\alpha_i\le Dd/d_i.
\]
Their weights are at most $sKDd$ after increasing the uniform
$K$. If $|m_{R,i}|>Dd/d_i$ for some $i$, the boxes have disjoint
exponent coordinates in that direction. Within each box the
independence of $b$ gives distinct values. There would be at
least $D^{m+1}\mathcal D_m$ distinct values. The same
cutoff-inclusion argument applies to the regular eigenfunction
products $R^jf^\alpha$, and bounds their number by
$C_WK^mD^m\mathcal D_m$. Choose $D>C_WK^m$,
fix this cutoff and all data, and then let $p$ be arbitrarily
large. It follows that
$|(\phi g_s-e)_i|\le(D+1)d/d_i$, with harmless adjustment of
the constant for the $o(1)$ term. Divide by the lower bound
for $s$ in \eqref{round:s-bounds} to obtain \eqref{round:U3}.

On a fixed branch, $A_\eta=e\cdot b$ and $s=g_s\cdot b$ are
rational-linear formulas. The derivative of their ratio is
$((e-\phi g_s)/s)\cdot\dot b$. At independent parameters the
coefficient vectors are exactly those of the selected vertex,
so \eqref{round:derivative-bound} follows.
\end{proof}

\subsection{Uniform constants and dependent endpoints}

We explain the order of constants before constructing the rounded
tuples. At most $n$ functions can be adjoined and at most $n$
tier rounds occur. There are at most a fixed multiple of $n$
approximate continuations. Reserve one unit of ratio increase
for every possible continuation and every tier round, together
with the one-unit margin for the stopping argument below.
Choose a single numerical bound $M_*$ larger than the initial
ratio one plus this entire reserve.

Upper weight bounds are fixed by a finite numerical recursion.
Start, for example, with $B_0=2$. At each possible enlargement,
apply the already proved estimate $C_5(B_j,M_*)$, multiply it
by the uniform tier-comparability bound for $d/d_{r+1}$, and
choose $B_{j+1}$ above that number and $B_j$. Add a further
fixed movement allowance at each possible small move or round.
The lower bound $d_i\ge\delta$ makes a bounded absolute
coordinate movement fit this allowance. This recursion has
predetermined finite depth; its definition does not assume
that any future tuple has already been constructed. Evaluate
the derivative constants and cutoff multipliers for all these
prospective bounds and retain their maxima. Only after these
choices will we select a large uniform positive integer $S$.
Finally omit a finite prefix of the sequence to put all the
finitely many fixed cutoffs below the next tier.

For the chosen element $\gamma\in G$, let
$\vartheta_f=\phaseangle\bigl(\chi_f(\gamma)\bigr)\in\R/\Z$. Targets will satisfy
\begin{equation}\label{round:targets}
 S b_j^*\equiv\vartheta_{f_j}\pmod\Z.
\end{equation}
For each germ use the movement tolerance
\begin{equation}\label{round:epsilon}
 \varepsilon=\frac{d_1}{Sd_n}
\end{equation}
in every coordinate of an approximate continuation; those
movements can always be chosen arbitrarily smaller.

Here is the construction at a tier endpoint $m$.
For the first tier start with its full original prefix $f_j^0$.
The valuation $v$ is Gauss with $b_j=d_j$, $s=1$, and
$\Phi=1$. Continue to very close independent weights, allowing
one unit of ratio increase. At a later tier start from the
previous completed tuple through $m_-$ at its exact target,
with the Gauss extension produced by the previous round, and
again continue to independent weights. While $r<m$, use
Lemma~\ref{round:enlarge} and then make a continuation within
the tolerance \eqref{round:epsilon}. Each step uses the
prospectively fixed bounds.

At $r=m$, every old coordinate $j\le m_-$ is within
$(n+1)\varepsilon$ of its previous target. Keep exactly that
old target, and choose for each new coordinate the nearest
target satisfying \eqref{round:targets}. Let $b^*$ be the
result and move along the straight segment from $b$ to $b^*$.
Its weighted length obeys
\begin{equation}\label{round:movement}
 \sum_{i\le m}\frac d{d_i}|b_i^*-b_i|\le\frac{C_7}{S}.
\end{equation}
For old coordinates this follows from
$d\le d_n$, $d_1\le d_i$, and \eqref{round:epsilon}; there
are at most $n$ of them. For new coordinates, the distance
to the nearest target is at most $1/(2S)$ and $d/d_i$ is
uniformly bounded within their tier.

All required positivity and anchor conditions hold for a
sufficiently large common $S$. The anchor $f_1=f_1^0$ begins
arbitrarily close to $d_1$; after its first round its exact
target never changes, and subsequent accumulated errors before
rounding are bounded as above. Thus
$1/2\le b_1/d_1\le2$ throughout. Every newly selected function
is nonconstant, and before rounding its coordinate is at least
$s v(f_i)\ge c d_1$ by \eqref{val:comparison}. The uniform lower
bound for $d_1$ keeps all targets and segments positive for
large $S$. The coordinate movements are $O(n/S)$, so the
prospective upper bounds also remain valid.

\begin{lemma}\label{round:segment}
For the bounds fixed above, choose $S$ large enough that
$C_6C_7/S<1$. If a completed prefix at the start of the segment
has a Gauss extension with an available ratio bound $M$, its
target endpoint has a centred quasi-monomial Gauss extension
with ratio at most $M+1$.
\end{lemma}

\begin{proof}
A constant segment requires no change. Otherwise independent
parameters are dense along the segment: for each nonzero
rational vector a rational dependence is an affine equation
in the segment parameter, not identically zero because the
starting weights are independent. Apply the finite interval
partition of Proposition~\ref{val:optimization}. Initially
approximate continuation supplies feasible independent
assignments with limiting optimal ratio at most $M$.

On each open interval where a selected branch is feasible,
integrate \eqref{round:derivative-bound} as long as the ratio
stays within the reserved bound $M+1$. The estimate holds at
the dense independent parameters and hence everywhere on
the interval by smoothness of the rational branch formula.
At a terminal parameter of such a bounded interval, selected
vertices at independent parameters have a compact convergent
subsequence on the fixed complex: $s\le2$, $A_\eta$ is bounded,
discrepancy is proper on the relevant cones, and
\eqref{round:s-bounds} keeps $s$ positive. Its limit is Gauss
on the tuple by Proposition~\ref{val:optimization}, including
all polynomial identities at these possibly dependent weights.
Lemma~\ref{val:gauss-tools} continues this valuation to nearby
independent assignments to the right with limiting ratio no
larger than the endpoint value. Thus there is neither loss
of feasibility nor an upward jump at a branch endpoint.

Proceed in order through the finitely many intervals of the
partition. On an interval, stop at a putative first excess
of the reserved budget; until that time all constants are
valid. With the segment parameterized by $0\le t\le1$,
integration and \eqref{round:movement} propagate the bound
\[
 \Phi\le M+t\frac{C_6C_7}{S}.
\]
The limsup argument propagates it across every endpoint.
Since $C_6C_7/S<1$, a first excess of $M+1$ is impossible.
The polyhedral selection may be relaxed at dependent
parameters inside intervals; this causes no difficulty,
because the selected branch is actual at the dense
independent parameters where the derivative is tested, and
compact limits provide actual Gauss extensions at endpoints.
At the final endpoint take the compact limit itself.
\end{proof}

The preceding numerical recursion and Lemma~\ref{round:segment}
justify the construction inductively through every tier.
There is no circular dependence in the choice of $S$:
all prospective upper bounds, ratio budgets, derivative
constants, movement constants, and cutoff multipliers precede
it. Increasing $S$ once enforces every one of the finitely
many conditions. By the final tier we obtain $n$ centred
regular eigenfunctions, positive rational exact targets
$b_i^*$, and a centred quasi-monomial Gauss extension $w$
such that
\begin{equation}\label{round:final-bounds}
 S b_i^*\equiv\vartheta_{f_i}\pmod\Z,\qquad
 s(w)\ge c,\qquad A_\eta(w)\le C,\qquad
 b_i^*\le B_*d_i.
\end{equation}
The target weights are rational because the character angles
are rational. A Gauss valuation on a full tuple with positive
rational weights is divisorial-scaled: its value group has
rational rank one and its residue field has transcendence
degree $n-1$. Finite field extension preserves these
properties, or one can use the smooth-stratum description
of Lemma~\ref{val:gauss-tools}. Thus the resulting $w$ is
divisorial-scaled.

\subsection{All rational semi-invariants}

\begin{proof}[Completion of the proof of Theorem~\ref{thm:phase}]
It remains to pass from the congruences on the completed tuple
to every rational semi-invariant, with one fixed additional
integer. Choose a uniform positive integer $D_{\rm fin}$ so
large that
\begin{equation}\label{round:final-count-constant}
 (D_{\rm fin}+1)\left(\frac{c}{2nB_*}\right)^n>C_W.
\end{equation}
For any regular nonzero eigenfunction $h$, take $T\ge d_n$
sufficiently large and consider the indexed functions
\begin{equation}\label{round:all-character-box}
 h^jf^\alpha,\qquad 0\le j\le D_{\rm fin},\qquad
 0\le\alpha_i\le\frac{sT}{2n b_i^*},\quad\alpha_i\in\Z.
\end{equation}
The number of indices is at least
\[
 (D_{\rm fin}+1)\left(\frac{c}{2nB_*}\right)^n
                  \prod_{i=1}^n\frac T{d_i}
 >\mathcal H(T).
\]
Choose $T$ still larger so that $D_{\rm fin}w(h)\le sT/2$.
Every indexed function then has $w$-value at most $sT$.
There is a nontrivial linear dependence modulo $(v>T)$,
whether or not some of the indexed functions already coincide.
Since $v$ is invariant, the ideal is stable and the dependence
can be chosen within a single character. Its sum either
vanishes or has $w$-value greater than $sT$, by
\eqref{val:comparison}. Hence the lowest terms must cancel
in the $w$-graded ring.

This cancellation involves at least two distinct exponents
$j,j'$. Otherwise, after factoring the nonzero initial of
$h^j$, a nonzero polynomial in the initials of $f$ would
vanish. This contradicts their Gauss property, even though
the target weights have rational relations. Choose two
lowest terms with distinct powers. Equality of their weights
and equality of their characters, together with
\eqref{round:final-bounds}, give
\begin{equation}\label{round:torsion}
 (j-j')\bigl(Sw(h)-\vartheta_h\bigr)=0
                      \quad\text{in }\R/\Z,
 \qquad 1\le|j-j'|\le D_{\rm fin}.
\end{equation}
Set
\begin{equation}\label{round:phase-valuation}
 \ell=\operatorname{lcm}(1,\ldots,D_{\rm fin}),\qquad
 w'=\ell S w.
\end{equation}
Then for every regular nonzero eigenfunction,
\[
 w'(h)\equiv\ell\,\phaseangle\bigl(\chi_h(\gamma)\bigr)\pmod\Z.
\]
The cutoff $T$ may depend on $h$ and on the germ; the integers
$D_{\rm fin},\ell,S$ do not. In particular the same $\ell$
works for all these functions at once.

For a rational semi-invariant $h$, choose a regular denominator
$b$ and replace it by the product of its group translates.
This is a nonzero invariant regular denominator $b_G$, and
$hb_G$ is regular and has the same character as $h$.
Subtracting the two already established congruences proves
the assertion for $h$. On invariant rational functions the
resulting value is integral. Finally $w'$ is centred and
divisorial-scaled, and its discrepancy is bounded by
$\ell S C$ uniformly on the retained subsequence. This is
exactly Theorem~\ref{thm:phase}.

The only input whose proof has been deferred is
Proposition~\ref{prop:integral-jump}, used in
Proposition~\ref{prop:product}. The remaining sections prove
that proposition and thereby complete the dependency chain.
\end{proof}
